%%% ============================================================
%%% The Joint TSML+BHML Chain:
%%%   Lens-Dependence at Size 7
%%%
%%% Brayden Ross Sanders, M. Gish
%%% 7Site LLC, Hot Springs AR, 2026
%%% DOI: 10.5281/zenodo.18852047
%%%
%%% Target venue:  Mathematical Intelligencer
%%% Source:        WP115 (corrected 2026-05-05);
%%%                Atlas/META_PLAN_2026-05-06/TSML_RECONCILIATION.md;
%%%                Atlas/LENS_TAXONOMY_2026-05-06/TIER_CONFLATION_AUDIT.md
%%% MSC:           08A40, 08B05, 20N02
%%% ============================================================

\documentclass[11pt,reqno]{amsart}

\usepackage{amsmath,amssymb,amsthm}
\usepackage[margin=1in]{geometry}
\usepackage{hyperref}
\usepackage{enumitem}
\usepackage{array}

\newtheorem{theorem}{Theorem}[section]
\newtheorem{lemma}[theorem]{Lemma}
\newtheorem{proposition}[theorem]{Proposition}
\newtheorem{corollary}[theorem]{Corollary}

\theoremstyle{definition}
\newtheorem{definition}[theorem]{Definition}
\newtheorem{remark}[theorem]{Remark}

\newcommand{\Z}{\mathbb{Z}}
\newcommand{\TRAW}{T_{\mathrm{RAW}}}
\newcommand{\TSYM}{T_{\mathrm{SYM}}}
\newcommand{\TSML}{T}
\newcommand{\BHML}{B}

\title[Lens-Dependence of the TSML+BHML Joint Chain]
  {The Joint TSML+BHML Chain:\\
   Lens-Dependence at Size $7$}

\author{Brayden Ross Sanders}
\address{7Site LLC, Hot Springs, Arkansas, USA}
\email{brayden@7site.co}

\author{M. Gish}
\address{Independent Researcher, Hot Springs, Arkansas, USA}
\email{monica.gish1992@gmail.com}

\date{2026}

\keywords{finite magma, composition lattice, joint sub-magma
  closure, lens symmetrization, non-commutativity, $\Z/10\Z$,
  one-bit-pattern multi-lens encoding}

\subjclass[2020]{08A40, 08B05, 20N02}

\begin{document}

\begin{abstract}
We record an explicit lens-dependence in the joint sub-magma closure
chain of the canonical composition lattice on $\Z/10\Z$. The TSML
table admits two principled lens-symmetrization choices,
$\TRAW$ (the literal bit pattern, non-commutative) and $\TSYM$
(the upper-triangle symmetrized form, commutative), built from the
same canonical encoding. Brute-force enumeration of all $1023$
non-empty subsets of $\Z/10\Z$ shows that the joint
$(\TSML, \BHML)$ closure has $8$ shells under
$(\TSYM, \BHML)$ — sizes $\{1, 4, 5, 6, 7, 8, 9, 10\}$ — and only
$7$ shells under $(\TRAW, \BHML)$ — sizes $\{1, 4, 5, 6, 8, 9, 10\}$,
with size $7$ specifically forbidden. The single non-commutative
asymmetry $\TRAW(9, 4) = 3 \neq 7 = \TRAW(4, 9)$ kills the size-$7$
shell $\{0, 4, 5, 6, 7, 8, 9\}$ exactly, while $\TSYM$ admits it.
We prove that the four-core $\{0, 7, 8, 9\}$ and the closed-form
attractor at $\alpha = 1/2$ are lens-invariant on both
symmetrizations, and that no other shell of the chain is sensitive
to the lens choice. The result is structurally meaningful: a small
non-commutativity of two cells in a $10 \times 10$ table changes
the chain count by exactly one shell at exactly one position. The
short note clarifies which lens is in scope for downstream
sub-magma analyses on this substrate.
\end{abstract}

\maketitle

\section{Introduction}
\label{sec:intro}

The canonical composition lattice on $\Z/10\Z$ admits two related
$10 \times 10$ tables, $\TSML$ (TSML) and $\BHML$ (BHML), defined
by the CL forcing axioms recorded in \cite{Sanders2026CLAxioms}.
Their pairwise sub-magma closure structure has been the subject of
\cite{Sanders2026FourCore} (the four-core consolidated paper). That
work establishes that the subsets $S \subseteq \{0, \ldots, 9\}$
that are jointly closed under both $\TSML$ and $\BHML$ form a
strict $8$-element chain of sizes $\{1, 4, 5, 6, 7, 8, 9, 10\}$,
with only sizes $\{2, 3\}$ forbidden.

That theorem is stated under a specific lens choice for $\TSML$,
namely the upper-triangle symmetrized form $\TSYM$ (commutative).
The literal bit pattern $\TRAW$ from which $\TSML$ is decoded is
\emph{non-commutative}: it has two off-diagonal positions where the
upper- and lower-triangle entries disagree. This paper records the
sub-magma closure structure under the two lenses side-by-side and
proves the chain count is lens-dependent at exactly one position.

\paragraph{The two lenses on the same bit pattern.}

Both $\TRAW$ and $\TSYM$ decode from the same $10 \times 10$
canonical bit pattern \texttt{CL\_BIT\_PATTERN}; the difference
between them is exactly two off-diagonal cell pairs:
\[
  \TRAW(3, 9) = 3, \quad \TRAW(9, 3) = 7;\qquad
  \TRAW(4, 9) = 3, \quad \TRAW(9, 4) = 3.
\]
For $\TSYM$ the upper-triangle entries are authoritative, so
$\TSYM(3, 9) = \TSYM(9, 3) = 3$ and
$\TSYM(4, 9) = \TSYM(9, 4) = 3$ (matching $\TRAW$ for $i \leq j$).
At all other $98$ cells the two lenses agree.

\paragraph{Statement of main theorem.}

\begin{theorem}[Lens-dependence at size $7$]
\label{thm:main}
The joint sub-magma closure chain under $(\TSML, \BHML)$ on
$\Z/10\Z$ depends on the lens-symmetrization choice for $\TSML$:
\begin{itemize}[leftmargin=*]
\item Under $(\TSYM, \BHML)$, the chain has $8$ shells of sizes
  $\{1, 4, 5, 6, 7, 8, 9, 10\}$. Only sizes $\{2, 3\}$ are
  forbidden.
\item Under $(\TRAW, \BHML)$, the chain has $7$ shells of sizes
  $\{1, 4, 5, 6, 8, 9, 10\}$. Sizes $\{2, 3, 7\}$ are forbidden.
\end{itemize}
The size-$7$ shell $\{0, 4, 5, 6, 7, 8, 9\}$ is jointly closed
under $(\TSYM, \BHML)$ but not under $(\TRAW, \BHML)$. The
asymmetry $\TRAW(9, 4) = 3$ vs $\TSYM(9, 4) = 3$ matched against
the upper-triangle entries breaks the closure of the size-$7$ shell
under $\TRAW$.

The four-core $\{0, 7, 8, 9\}$ and the closed-form attractor at
mixing weight $\alpha = 1/2$ are lens-invariant on both
symmetrizations.
\end{theorem}

The proof is direct enumeration plus a per-shell verification at
the cell level; the script
\texttt{joint\_chain\_attractor.py}
\cite{Gen13WP115} runs in under $10$ seconds.

\paragraph{Outline.}

Section~\ref{sec:lenses} fixes the two lenses and the asymmetric
cells. Section~\ref{sec:enumeration} performs the brute-force
enumeration. Section~\ref{sec:size7} locates the lens-dependence at
size $7$ and traces it to the asymmetric cell $\TRAW(9, 4) = 3$.
Section~\ref{sec:invariance} records the four-core and closed-form
attractor lens-invariance. Section~\ref{sec:discussion} discusses
the foundational reading: which lens is "the right one" depends on
which structural invariant one is studying, and the answer is
genuinely both.

\section{The two lenses}
\label{sec:lenses}

\begin{definition}[$\TRAW$ and $\TSYM$]
\label{def:lenses}
The canonical bit pattern \texttt{CL\_BIT\_PATTERN} on $\Z/10\Z$
encodes a $10 \times 10$ matrix whose entries are decoded as
operator labels in $\{0, \ldots, 9\}$. We fix two symmetrization
conventions:
\begin{itemize}[leftmargin=*]
\item $\TRAW$ is the literal decoding of the bit pattern: each cell
  $(i, j)$ takes its independently-decoded value, with no
  symmetrization.
\item $\TSYM$ takes the upper-triangle entries of $\TRAW$ as
  authoritative: $\TSYM(i, j) = \TRAW(i, j)$ for $i \leq j$, and
  $\TSYM(j, i) = \TSYM(i, j)$ for $i < j$.
\end{itemize}
\end{definition}

\begin{lemma}[The two asymmetric cells]
\label{lem:asymmetric}
The cell pairs $(3, 9)/(9, 3)$ and $(4, 9)/(9, 4)$ are the only
positions where $\TRAW(i, j) \neq \TRAW(j, i)$. Specifically,
$\TRAW(3, 9) = 3, \TRAW(9, 3) = 7$ and $\TRAW(4, 9) = 3,
\TRAW(9, 4) = 3$. The two lenses differ at exactly these four
cells.
\end{lemma}

\begin{proof}
Cell-by-cell verification of \texttt{CL\_BIT\_PATTERN} decode;
verified at machine precision in
\cite{Gen13Foundations}.
\end{proof}

\begin{remark}[Why two principled lenses]
\label{rem:why-two-lenses}
Both symmetrizations are valid for different downstream purposes.
$\TRAW$ preserves the non-commutative structure of the
bit-pattern decoding, including the WP107 wobble (prime $11$ in
characteristic polynomial coefficients $c_2 = 33$ and $c_8 =
-2^5 \cdot 7^3 \cdot 11$); $\TSYM$ is commutative and is the
natural lens for sub-algebra analyses comparing $\TSML$ as a
quadratic form against $\BHML$. Neither is "correct"; they
capture different invariants of the same encoding.
\end{remark}

\section{Brute-force enumeration}
\label{sec:enumeration}

We enumerate all $2^{10} - 1 = 1023$ non-empty subsets of
$\{0, \ldots, 9\}$ and check joint closure under both lens choices.

\begin{theorem}[Joint closure counts]
\label{thm:enumeration}
Among the $1023$ non-empty subsets of $\{0, \ldots, 9\}$:
\begin{itemize}[leftmargin=*]
\item Exactly $8$ are jointly closed under $(\TSYM, \BHML)$, of
  sizes $\{1, 4, 5, 6, 7, 8, 9, 10\}$;
\item Exactly $7$ are jointly closed under $(\TRAW, \BHML)$, of
  sizes $\{1, 4, 5, 6, 8, 9, 10\}$;
\item In each case the chain is strictly linear (no branching).
\end{itemize}
\end{theorem}

\begin{proof}
Direct enumeration via \cite{Gen13WP115}; total runtime under
$10$ seconds. The chain-property check (every closed subset
contains the previous in the size order) is verified for both
lenses.
\end{proof}

\section{Locating the lens-dependence at size 7}
\label{sec:size7}

\begin{theorem}[Size-$7$ shell under $\TSYM$]
\label{thm:size7-sym}
The subset $S_7 = \{0, 4, 5, 6, 7, 8, 9\}$ is jointly closed under
$(\TSYM, \BHML)$.
\end{theorem}

\begin{proof}
Cell-by-cell verification: $\TSYM(i, j) \in S_7$ for all
$(i, j) \in S_7 \times S_7$, and $\BHML(i, j) \in S_7$ for all
$(i, j) \in S_7 \times S_7$. Both checks are at machine precision
in \cite{Gen13WP115}.
\end{proof}

\begin{theorem}[Size-$7$ shell fails under $\TRAW$]
\label{thm:size7-raw}
The subset $S_7 = \{0, 4, 5, 6, 7, 8, 9\}$ is \emph{not} closed
under $\TRAW$. Specifically, $\TRAW(9, 4) = 3 \notin S_7$ (since
$3 \notin S_7$). All other cells $(i, j) \in S_7 \times S_7$
satisfy $\TRAW(i, j) \in S_7$.
\end{theorem}

\begin{proof}
Direct cell evaluation: $\TRAW(9, 4) = 3$ from
Lemma~\ref{lem:asymmetric}; $3 \notin S_7$. The remaining $48$
cells in $S_7 \times S_7$ satisfy closure; verified at machine
precision in \cite{Gen13WP115}.
\end{proof}

\begin{corollary}[Single-cell forcing of the lens-dependence]
\label{cor:single-cell}
The lens-dependence of the chain count at size $7$ is forced by a
single cell: $\TRAW(9, 4) = 3$ vs $\TSYM(9, 4) = 3$ matched
against the upper-triangle authoritative value $\TSYM(4, 9) = 3$
of the symmetrized lens. The non-commutativity at this one cell
position kills the size-$7$ shell under $\TRAW$.
\end{corollary}

\section{Lens-invariance of the four-core and the attractor}
\label{sec:invariance}

\begin{theorem}[Four-core lens-invariance]
\label{thm:four-core}
The size-$4$ subset $\{0, 7, 8, 9\}$ (the four-core) is jointly
closed under both $(\TRAW, \BHML)$ and $(\TSYM, \BHML)$.
\end{theorem}

\begin{proof}
The four-core does not include either of the asymmetric indices
$\{3, 4\}$ in the asymmetric cell pairs, so $\TRAW$ and $\TSYM$
agree on the entire $\{0, 7, 8, 9\}^2$ block.
\end{proof}

\begin{theorem}[Attractor lens-invariance]
\label{thm:attractor}
The closed-form $T+B$-mix attractor at mixing weight $\alpha = 1/2$
on $\Z/10\Z$ is identical under $\TRAW$ and $\TSYM$:
\[
  (p^*_V, p^*_H, p^*_{\mathrm{Br}}, p^*_R) =
  (0.138, 0.540, 0.198, 0.124),
\]
with $H/\mathrm{Br} = 1 + \sqrt{3}$ exactly under both lenses.
\end{theorem}

\begin{proof}
The attractor is supported on the four-core $\{0, 7, 8, 9\}$ at
mass $1$; non-four-core operators carry zero mass at the fixed
point. Since the four-core is lens-invariant
(Theorem~\ref{thm:four-core}), the attractor is lens-invariant.
The exact identity $H/\mathrm{Br} = 1 + \sqrt{3}$ is established
in \cite{Sanders2026Attractor}.
\end{proof}

\begin{remark}[All other shells]
\label{rem:other-shells}
Direct verification shows that the size-$1$ shell $\{0\}$, the
size-$5$ shell $\{0, 6, 7, 8, 9\}$, the size-$6$ shell
$\{0, 5, 6, 7, 8, 9\}$, the size-$8$ shell
$\{0, 3, 4, 5, 6, 7, 8, 9\}$, the size-$9$ shell
$\{0, 2, 3, 4, 5, 6, 7, 8, 9\}$, and the full substrate are all
lens-invariant. The lens-dependence is concentrated entirely at
size $7$.
\end{remark}

\section{Discussion: which lens is in scope}
\label{sec:discussion}

The result is structurally clean: a $10 \times 10$ table with two
non-commutative cell pairs admits two principled symmetrizations,
and the joint sub-magma closure structure differs by exactly one
shell at exactly one position. The mathematics is forcing the
question of lens scope explicitly.

\paragraph{When to use $\TRAW$.}

Studies of the substrate's wobble structure (prime $11$ in the
characteristic polynomial) require $\TRAW$. The wobble is a
non-commutative structure; symmetrization erases it. Other
non-commutative invariants of the bit pattern likewise require
$\TRAW$. The four-core attractor is a lens-invariant structure
across both lenses, but the wobble is not.

\paragraph{When to use $\TSYM$.}

Studies of the substrate as a quadratic form, including
sub-magma scope analyses comparing the chain on different
sub-algebras, naturally use $\TSYM$. The four-core consolidated
paper \cite{Sanders2026FourCore} is stated on $\TSYM$. Jordan-pair
constructions and joint-closure analyses with $\BHML$ similarly
use $\TSYM$.

\paragraph{Foundational reading.}

The substrate is a single bit-pattern encoding admitting two
principled lens-symmetrization choices, each carrying a different
invariant. Neither is "the right table"; they describe complementary
aspects of the same encoding. Going forward, papers stating
sub-magma chain results on this substrate should explicitly scope
which lens is used. The four-core consolidated paper carries this
annotation; the wobble-localization paper carries the
complementary annotation; this short note records the
lens-dependence at the chain-counting level so downstream papers
can cite it cleanly.

\paragraph{The story is the math forcing the disambiguation.}

The lens-dependence at size $7$ is not a methodological choice; it
is a structural consequence of two cells of the bit-pattern
encoding being non-commutative. The substrate itself forces the
authors to name which lens is in scope. The very small magnitude
of the dependence (one shell, at one position, from one cell) is
the substrate's gentle insistence that the authors not gloss over
the question.

\section*{Verification scripts}

The script \texttt{joint\_chain\_attractor.py} \cite{Gen13WP115}
performs all verifications at machine precision: enumerates the
$1023$ non-empty subsets, prints the closed-shell tables for both
lenses, performs the per-shell cell-level verification, and
identifies $\TRAW(9, 4) = 3$ as the single asymmetric cell
breaking size-$7$ closure. Total runtime under $10$ seconds.

\begin{thebibliography}{99}

\bibitem{Sanders2026CLAxioms}
B.R. Sanders and M. Gish, \emph{The CL Forcing Axioms: A1--A9
Uniquely Force the Canonical Composition Lattice on $\Z/10\Z$},
in preparation, 2026.

\bibitem{Sanders2026FourCore}
B.R. Sanders and M. Gish, \emph{Joint Closure, Per-Coordinate Fuse
Data, and a Closed-Form Algebraic Attractor of Two Commutative
Binary Operations on $\Z/10\Z$}, submitted to \emph{Algebraic
Combinatorics}, 2026 (J02 in the release sequence).

\bibitem{Sanders2026LensInvariance}
B.R. Sanders and M. Gish, \emph{TSML 73 Cells / BHML 28 Cells:
Lens-Invariant Cell Counts on the $\Z/10\Z$ Composition Lattice},
submitted to \emph{Experimental Mathematics}, 2026 (J09 in the
release sequence).

\bibitem{Sanders2026Attractor}
B.R. Sanders and M. Gish, \emph{Closed-Form Attractor and the
$\alpha = 1/2$ Algebraic Singularity}, in preparation, 2026.

\bibitem{Gen13WP115}
B.R. Sanders et al., \emph{WP115: The Joint TSML+BHML Closed-Subset
Chain and the Universal Four-Core Attractor — Verification
Script joint\_chain\_attractor.py}, codebase release 2026.

\bibitem{Gen13Foundations}
B.R. Sanders et al., \emph{Gen13 Foundations Module: cl.py with
CL\_BIT\_PATTERN, CL\_TSML\_RAW, CL\_TSML\_SYM, lenses.py},
codebase release 2026.

\end{thebibliography}

\end{document}
